%---------------------------- CHAPTER ----------------------------
\chapter{Demolition and Refurbishment Risk Assessment}
\label{chap:Demolition and Refurbishment Risk Assessment}

%---------------------------- SECTION ----------------------------
\section{Why this assessment matters in construction}

\paragraph{}
Demolition and refurbishment operations represent some of the most hazardous activities undertaken in the UK construction industry. Unlike new-build projects, where hazards are introduced progressively as the structure rises, demolition presents the inverse challenge: the structure being removed may be unstable, contaminated, or concealing hidden dangers that are not apparent until work is already under way. The Health and Safety Executive's injury statistics consistently identify demolition as one of the highest-risk trade sectors in construction, with fatal and major injury rates significantly above the industry average. The combination of structural unpredictability, hazardous materials — most notably asbestos-containing materials (ACMs) — connected services, dust, noise, vibration, and the compressed timescales typically imposed by clients creates a hazard profile that demands meticulous pre-work planning and a rigorously documented risk assessment process.

\paragraph{}
Refurbishment projects introduce many of the same hazards as full demolition but compound them with the presence of an occupied or partially occupied built environment. Stripping out internal finishes, removing partition walls, cutting through floor slabs, and altering load-bearing elements all carry the potential for sudden structural movement, release of fibrous or chemical contaminants, and interference with live services. The risk assessment process for refurbishment must therefore consider not only the tradespeople undertaking the work but also any building occupants, neighbouring properties, and members of the public who may be in proximity. Effective planning begins well before the first operative sets foot on site, requiring a structured survey and information-gathering phase that informs every subsequent decision about method, sequence, and control.

\paragraph{}
The legislative framework governing demolition and refurbishment in the United Kingdom is extensive. The Construction (Design and Management) Regulations 2015 (CDM 2015) place specific duties on clients, principal designers, and principal contractors to manage pre-construction information, produce and maintain a Construction Phase Plan, and ensure that the health and safety file is updated to reflect any changes to the built asset. BS 6187:2011 — the Code of Practice for Full and Partial Demolition — provides the industry benchmark for planning and executing demolition works, covering everything from preliminary structural surveys to the selection of demolition methods and the management of arisings. Together, these instruments create a framework within which the risk assessment must be anchored, ensuring that professional judgement is exercised systematically and that decisions can be justified by reference to recognised good practice.

\paragraph{}
The consequences of inadequate risk assessment in demolition are severe and well-documented. Structural collapses during demolition have caused multiple fatalities at sites across the UK, frequently attributable to failure to commission an adequate pre-demolition structural survey, failure to identify and isolate connected services before work commenced, or failure to recognise and protect adjacent structures from the effects of demolition activity. The economic consequences — project delays, prosecutions, civil claims, reputational damage — are equally significant. A thorough risk assessment, conducted by a competent person and integrated into a site-specific method statement that has been reviewed and approved by a structural engineer where stability is at issue, is not merely a regulatory obligation; it is the foundation upon which the safe execution of every demolition and refurbishment project depends.

\barquote{``Demolition is not the reverse of construction — it is a distinct discipline requiring its own body of knowledge, its own specialist competence, and its own rigorous approach to risk. The structure being taken down has a history that must be understood before a single stone is moved.'' — National Federation of Demolition Contractors guidance}

\example{Demolition of a 1960s Office Block with Suspected ACMs}{A principal contractor was appointed to demolish a five-storey reinforced concrete office building constructed in the early 1960s. Before any physical works commenced, a pre-demolition refurbishment and demolition (R\&D) asbestos survey (see Chapter~\ref{chap:Asbestos Risk Assessment}) was commissioned in accordance with HSG264, identifying chrysotile-containing floor tiles, amosite insulation board in ceiling voids, and crocidolite pipe lagging in the basement plant room. The structural engineer reviewed the original construction drawings — obtained from the local authority building control archive — and identified that the building had been subject to an undisclosed basement extension in the 1980s that had partially removed a ground-floor transfer beam. The pre-demolition survey flagged this as a critical stability issue. The method statement and demolition sequence were revised accordingly, the asbestos was removed under licensed conditions prior to structural demolition commencing, and the basement was shored before the superstructure was progressively demolished by high-reach excavator from the roof downward. The risk assessment documented each of these decisions and their evidential basis, providing a clear audit trail for the principal contractor's CDM obligations.}

%---------------------------- SECTION ----------------------------
\section{Legal and professional framework}

The legal and professional framework governing demolition and refurbishment risk assessment in the United Kingdom is multi-layered, drawing on primary health and safety legislation, subordinate regulations, approved codes of practice, and technical standards published by the British Standards Institution. At the apex of the legislative hierarchy sits the Health and Safety at Work etc.\ Act 1974, which imposes a general duty on employers to ensure, so far as is reasonably practicable, the health, safety, and welfare of their employees and others who may be affected by their undertakings. This duty is given operational content by the Management of Health and Safety at Work Regulations 1999, which require employers to conduct suitable and sufficient risk assessments and to implement appropriate preventive and protective measures.

Within the construction sector specifically, CDM 2015 creates a structured project management framework that is particularly relevant to demolition and refurbishment. The regulations require that a principal contractor is appointed on projects involving more than one contractor, that a Construction Phase Plan is prepared before physical work commences, and that pre-construction information — including information about hazardous materials, structural history, and connected services — is assembled and communicated to all duty-holders. Notification to the HSE under CDM 2015 is required for demolition projects that will last longer than 30 working days with more than 20 workers simultaneously, or that will exceed 500 person-days. In practice, most significant demolition contracts will meet these thresholds and must be notified via the HSE's online F10 portal.

\begin{itemize}
\bitem{BS 6187:2011 — Code of Practice for Full and Partial Demolition}{The principal British Standard governing demolition methodology. It covers pre-demolition surveys, structural appraisal, method selection (mechanical, implosion, hand, saw-cutting), sequence of works, protection of adjacent structures, management of arisings, and safety management. It is the definitive technical reference for demolition method statements and risk assessments in the UK.}
\bitem{CDM 2015 — Construction (Design and Management) Regulations 2015}{Requires pre-construction information to be gathered and communicated, a Construction Phase Plan to be prepared and maintained, and — for notifiable projects — an F10 notification to be submitted to the HSE. Places specific duties on the principal contractor regarding on-site management of demolition and refurbishment risks.}
\bitem{Control of Asbestos Regulations 2012 (CAR 2012)}{Requires that a refurbishment and demolition (R\&D) asbestos survey is conducted by a UKAS-accredited surveyor before any demolition or intrusive refurbishment work begins, and that all identified ACMs are managed or removed in accordance with the regulations before structural demolition proceeds. See Chapter~\ref{chap:Asbestos Risk Assessment}.}
\bitem{HSG151 — Protecting the Public: Your Next Steps}{HSE guidance on protecting members of the public and others not employed on site from the risks arising from construction, demolition, and refurbishment activities, including dust, falling debris, structural collapse, and traffic management around demolition sites.}
\bitem{Party Wall etc.\ Act 1996}{Where demolition works are adjacent to or involve a party wall, the Act requires formal notices to be served on adjoining owners and, where agreement cannot be reached, for a party wall award to be made by an agreed or appointed surveyor. The structural survey must identify all party wall and boundary conditions before demolition commences.}
\bitem{Control of Pollution Act 1974 and Environmental Protection Act 1990}{Regulate noise and dust emissions from demolition sites, requiring best practicable means to be employed. Local authority consent under Section 61 of the Control of Pollution Act may be sought proactively to agree noise and vibration limits and working hours with the local authority in advance of works commencing.}
\bitem{Lifting Operations and Lifting Equipment Regulations 1998 (LOLER)}{Apply to crane operations used during demolition, including the lifting and placing of demolished sections, the use of demolition grabs and shears, and the erection and dismantling of temporary works. All lifting operations must be planned by a competent person, supervised, and carried out in a safe manner.}
\bitem{Provision and Use of Work Equipment Regulations 1998 (PUWER)}{Govern the selection, inspection, maintenance, and use of demolition plant including high-reach excavators, hydraulic breakers, wire rope attachments, and saw-cutting equipment. Operators must be trained and competent, and equipment must be inspected before use.}
\end{itemize}

%---------------------------- SECTION ----------------------------
\section{Conducting the assessment using the five-step method}

\subsection{Step 1 --- Identify Hazards}

\paragraph{}
The hazard identification phase for demolition and refurbishment must begin with a thorough review of all available pre-construction information, including original construction drawings, structural calculations, as-built records, previous survey reports, planning history, building control records, and any known maintenance history. Where original drawings are unavailable — which is common for structures predating the 1970s — a structural engineer must undertake an intrusive investigation to establish the construction method, material properties, and load path before demolition planning can proceed. The pre-demolition structural survey required by BS 6187:2011 must identify all structural elements and their sequence of load transfer, the condition of all structural materials, any modifications made since original construction, the presence of hazardous materials (ACMs, lead paint, contaminated ground), the location and condition of all connected services (gas, electricity, water, telecommunications, district heating), and any geotechnical factors affecting stability during the demolition sequence. The R\&D asbestos survey — mandatory under CAR 2012 — must be completed before any intrusive structural work begins, and all licensed asbestos removal must be completed before structural demolition commences, except where specific risk-assessed working arrangements have been agreed with the licensed contractor.

\subsection{Step 2 --- Decide Who or What Is Affected}

\paragraph{}
The affected parties in a demolition and refurbishment project extend well beyond the operative wielding the demolition tool. Structural collapse, whether partial or total and whether planned or unplanned, can affect workers throughout the building or site, operatives on adjacent structures, members of the public on footways and roads adjacent to the site, occupants of neighbouring properties, and utility workers in the highway. Dust and noise generated by demolition operations (see Chapter~\ref{chap:Silica, Wood Dust and Hazardous Dust Risk Assessment} and Chapter~\ref{chap:Noise, Vibration and Hand-Arm Vibration Risk Assessment}) affect not only those within the exclusion zone but also those in the wider neighbourhood. Groundborne vibration from mechanical demolition can damage adjacent structures — particularly older masonry buildings — and cause distress and disturbance to neighbouring occupants. The risk assessment must identify each of these receptor groups and assess the pathways by which harm could reach them, so that control measures can be targeted proportionately.

\subsection{Step 3 --- Evaluate Likelihood and Consequence}

\paragraph{}
The likelihood and consequence evaluation for demolition hazards must consider both the inherent risk — prior to any controls — and the residual risk that remains after controls are implemented. Structural collapse during demolition carries a consequence rating of catastrophic, since even a partial collapse of a multi-storey structure can result in multiple fatalities. The likelihood of such an event is strongly influenced by the quality of the pre-demolition survey, the adequacy of the method statement and demolition sequence, and the competence of those executing the works. Where structural uncertainty exists — for example, where original drawings cannot be located or where intrusive investigation reveals unexpected structural modifications — the inherent likelihood must be rated as high until a structural engineer has reviewed the situation and confirmed a safe sequence. ACM release during demolition carries a high consequence rating for mesothelioma risk over a latency period of decades; the likelihood of release is directly related to the thoroughness of the R\&D survey and the quality of the licensed removal works.

\subsection{Step 4 --- Select and Implement Controls}

\paragraph{}
Controls for demolition and refurbishment hazards must be selected systematically, working down through the hierarchy of controls (see Section~\ref{sec:hierarchy} of this chapter). The starting point for structural hazards is engineering control: the demolition sequence must be designed so that the structure is never placed in a condition of instability, which requires the progressive and controlled removal of elements in the reverse order of their installation — most recently installed first, as required by BS 6187:2011. Where the structure cannot be made inherently stable throughout the demolition sequence, temporary works — shoring, propping, kentledge — must be designed and installed by a competent structural engineer. Exclusion zones must be established around the demolition face, sized by reference to the potential collapse envelope of the structure being demolished, and enforced by physical barriers, banksman attendance, and site access controls. Connected services must be isolated at source — not merely at the boundary of the building — before any demolition work commences, with written confirmation of isolation obtained from each utility provider or the client's building services engineer.

\subsection{Step 5 --- Record and Review}

\paragraph{}
The risk assessment and method statement (RAMS) for demolition must be documented in writing, reviewed and approved by a competent person — which for structural demolition means a structural engineer — and communicated to all operatives before work begins. The Construction Phase Plan must be updated to reflect the demolition RAMS and any significant changes to the demolition sequence or method that arise during the works. Daily briefings (toolbox talks) should be conducted to ensure that operatives are aware of the specific activities planned for each shift and the controls that apply. The risk assessment must be reviewed whenever there is a significant change in circumstances — discovery of unexpected structural conditions, identification of previously unrecorded ACMs, damage to the structure from weather or adjacent activity, or any near-miss or incident. At the completion of the project, lessons learned should be captured and fed back into the organisation's risk management procedures to improve future assessments.

\example{Five-Step Application --- Demolition of a Victorian Mill Building}{A demolition contractor was appointed to demolish a four-storey Victorian mill building constructed of load-bearing brick masonry with timber floors and a slate roof. Step 1 identified the following principal hazards: structural instability (timber floor joists rotted, roof partially collapsed), ACMs (chrysotile rope gaskets in boiler room, asbestos cement roofing on later extension), lead paint on structural ironwork, and a live gas main within 600 mm of the building's east wall. Step 2 identified affected parties as demolition operatives, a live footway on the south side of the building, and residential properties 8 m to the north. Step 3 rated structural collapse as catastrophic/medium likelihood (inherent), reduced to catastrophic/low following engineering controls. Step 4 controls included: structural engineer sign-off of sequence (east and north walls last, propped until completion), hoarding and scaffold fan to protect footway, gas isolated at governor by gas network operator, licensed asbestos removal of boiler room materials and full R\&D survey of extension before demolition, high-reach excavator working from south to north. Step 5 included daily RAMS briefings, structural engineer site visits at each key stage, and a recorded review following discovery of an unrecorded Victorian brick culvert beneath the east wall that required revision of the demolition sequence.}

%---------------------------- SECTION ----------------------------
\section{Applying the hierarchy of controls}

\paragraph{}
The hierarchy of controls must be applied rigorously in demolition and refurbishment risk assessment, recognising that the temptation to rely on personal protective equipment as a primary control is particularly acute in an industry where PPE is ubiquitous and its limitations are frequently underestimated. The hierarchy must be worked through systematically for each hazard identified, with the aim of eliminating or reducing risk at source before considering behavioural or equipment-based controls.

\begin{itemize}
\bitem{Elimination}{Remove the hazard entirely before demolition commences. All ACMs must be removed by a licensed contractor before structural demolition begins. All connected services must be isolated and disconnected at source. Where hazardous materials cannot be eliminated in advance — for example, lead paint encapsulated in structural metalwork — the demolition method must be designed to minimise disturbance and exposure.}
\bitem{Substitution}{Where a higher-risk demolition method can be replaced with a lower-risk alternative, substitution should be considered. Hand demolition of isolated structural elements may be safer than machine demolition in confined or structurally sensitive areas. Deconstruction — the selective removal and reuse of building components — reduces the volume of material requiring mechanical processing and may reduce dust and noise generation compared to bulk mechanical demolition.}
\bitem{Engineering controls}{The primary engineering controls in demolition are the demolition sequence and temporary works design. The sequence must ensure that the structure is never in an unstable condition during the works. Exclusion zones physically prevent access to the collapse envelope. Dust suppression (water misting, localised extraction) reduces airborne particulate at source. Acoustic hoarding and barriers reduce noise transmission to receptors. High-reach demolition machines allow structural elements to be brought down in controlled sections from a safe standoff distance. Suspended loads must be controlled by rigging to prevent uncontrolled movement.}
\bitem{Administrative controls}{Method statements and RAMS communicate the safe system of work. Permit-to-work systems control access to exclusion zones and manage interfaces with live services. Toolbox talks ensure operatives understand daily activities and associated hazards. Structural engineer sign-off at key stages provides an independent check on stability. Monitoring of adjacent structures for settlement or movement (crack monitoring gauges, precise levelling) provides early warning of ground movement or structural distress.}
\bitem{PPE}{PPE is the last resort and must not be used as a substitute for engineering or administrative controls. Minimum PPE for demolition operatives includes hard hat, safety boots with midsole protection, high-visibility vest, eye protection, and hearing protection (in zones exceeding 85 dB(A)). Where dust exposures are above the workplace exposure limit, respiratory protective equipment (RPE) — minimum FFP3 disposable half-mask or powered air-purifying respirator (PAPR) — must be worn. For licensed asbestos work, full-face respirator with P3 filter or equivalent is required in accordance with CAR 2012.}
\end{itemize}

%---------------------------- SECTION ----------------------------
\section{Cadence of assessment and review triggers}

\paragraph{}
The demolition risk assessment is not a static document produced once at the outset of a project and filed away. The dynamic and unpredictable nature of demolition — where structural conditions can change with every machine bite, and where the discovery of unexpected hazardous materials or structural conditions can fundamentally alter the safe method of work — requires that the risk assessment is treated as a living document, subject to continuous review and capable of rapid revision. The Construction Phase Plan must specify the process by which the risk assessment will be reviewed and updated, and all key duty-holders must be aware of their obligation to escalate information that triggers a review.

\begin{itemize}
\bitem{Discovery of unexpected structural conditions}{Where intrusive investigation or the progression of demolition reveals structural conditions not identified by the pre-demolition survey — undisclosed modifications, material deterioration, unexpected load paths — demolition must be paused and the structural engineer consulted before work resumes. A revised RAMS must be produced and briefed to all operatives.}
\bitem{Discovery of previously unidentified ACMs}{Where demolition reveals materials suspected of containing asbestos that were not identified by the R\&D survey, demolition must cease in the affected area immediately. The materials must be sampled and analysed by a UKAS-accredited laboratory, and, if ACMs are confirmed, licensed removal must be completed before demolition resumes. The risk assessment must be updated to reflect the revised scope.}
\bitem{Significant weather events}{High winds, heavy rainfall, or freezing conditions can alter the structural stability of a partially demolished structure. The risk assessment should specify weather thresholds at which demolition activity must cease — typically wind speeds above Beaufort Force 6 (approximately 22 mph mean wind speed) for high-reach machine demolition, and lower thresholds for hand demolition at height.}
\bitem{Near-miss or incident}{Any near-miss, dangerous occurrence, or injury arising from demolition activities must trigger an immediate review of the risk assessment to determine whether controls were adequate and whether the method of work needs to be revised. Dangerous occurrences reportable under RIDDOR 2013 — including unintended collapse or partial collapse of a building or structure — must be reported to the HSE and will typically trigger a formal investigation.}
\bitem{Change in scope or programme}{Where the client or principal contractor proposes a change in demolition scope, sequence, or programme — for example, accelerating a planned sequence, removing a planned temporary works stage, or adding an additional structure to the demolition scope — the risk assessment must be revised before the change is implemented.}
\bitem{CDM notification threshold reached}{Where a project that was not initially notifiable under CDM 2015 grows in scope to exceed the notification thresholds (30 working days with more than 20 workers simultaneously, or 500 person-days), the F10 notification must be submitted and the Construction Phase Plan reviewed to ensure it meets the requirements for notifiable projects.}
\end{itemize}

%---------------------------- SECTION ----------------------------
\section{Common failure modes}

\begin{itemize}
\bitem{Inadequate pre-demolition survey}{The most common root cause of serious demolition incidents is failure to commission, or to act on the findings of, an adequate pre-demolition structural survey. Surveys that rely solely on visual inspection without intrusive investigation are frequently unable to identify deteriorated structural elements, undisclosed modifications, or hidden load paths. The risk assessment produced on the basis of an inadequate survey will be fundamentally flawed.}
\bitem{Failure to isolate connected services}{Unexploded ordnance of the construction world, live services within a structure being demolished — particularly gas mains and high-voltage electrical cables — have caused catastrophic incidents when demolition plant has ruptured them. The risk assessment must require written confirmation of isolation from every utility provider before demolition commences, and the method statement must specify trial excavation or probing to verify service positions where plans are uncertain.}
\bitem{Failure to manage the sequence of demolition}{Demolition of structural elements in the wrong sequence — particularly removal of elements that are providing lateral restraint to adjacent walls, or removal of lower floors before upper floors have been removed — can cause progressive structural collapse. The method statement must sequence each element of the works, and site supervision must enforce the sequence strictly.}
\bitem{Inadequate exclusion zones}{Exclusion zones that are too small, poorly marked, or not enforced allow workers and members of the public to enter the collapse envelope. The collapse envelope for a multi-storey structure extends significantly beyond the perimeter of the building; BS 6187:2011 provides guidance on calculating the minimum exclusion zone for different structural forms and demolition methods.}
\bitem{Complacency on familiar sites}{Experienced demolition operatives may become complacent on projects that appear straightforward, shortcutting pre-work checks, bypassing permit-to-work systems, or proceeding without waiting for structural engineer sign-off. The risk assessment and site management system must be designed to make safe behaviour the path of least resistance, with clear escalation routes for operatives who identify changing conditions.}
\bitem{Inadequate dust and noise management}{Dust from demolition — including respirable crystalline silica (RCS) from concrete, brick, and mortar (see Chapter~\ref{chap:Silica, Wood Dust and Hazardous Dust Risk Assessment}) — can cause irreversible lung disease. Noise and vibration from mechanical demolition (see Chapter~\ref{chap:Noise, Vibration and Hand-Arm Vibration Risk Assessment}) can cause permanent hearing damage and damage to adjacent structures. These hazards are frequently undercontrolled in demolition RAMS, with RPE and hearing protection relied upon in place of engineering controls such as water suppression, localised extraction, and selection of lower-noise plant.}
\end{itemize}

\example{Failure Pattern --- Partial Collapse During Selective Demolition}{A refurbishment contractor was engaged to remove a single-storey rear extension from a Victorian terraced property. The structural survey identified the extension as constructed of single-skin brickwork with a flat felt roof and noted that the rear wall of the main house was load-bearing. No temporary works were specified. During demolition of the extension roof, the operative removed the timber wall plate that was bearing on the party wall between the extension and the neighbouring property. The neighbouring party wall, which was in poor condition and had been relying on the wall plate for lateral restraint, partially collapsed, damaging the neighbouring property and injuring a member of the public on the adjacent footway. Investigation revealed that the pre-demolition survey had been conducted visually only and had not identified the wall plate's structural function. The risk assessment had not required structural engineer sign-off for what was considered a ``minor'' demolition. The incident resulted in a Prohibition Notice, prosecution under HSWA 1974, and a substantial civil claim from the injured party.}

%---------------------------- CHAPTER SUMMARY ----------------------------
\newpage
\section{Chapter Summary}

\begin{table}[H]
\centering
\begin{tabularx}{\textwidth}{>{\raggedright\arraybackslash}p{3.2cm}X}
\toprule
\textbf{Assessment Element} & \textbf{Operational Requirement} \\
\midrule
Legal/Professional Basis & BS 6187:2011, CDM 2015, CAR 2012, HSWA 1974, Management Regulations 1999, Party Wall etc.\ Act 1996, LOLER, PUWER. \\
Method & Apply the full five-step process and document inherent/residual risk. Commission pre-demolition structural survey and R\&D asbestos survey before works commence. Obtain structural engineer sign-off on method statement and sequence. \\
Controls & Elimination of ACMs and live services before structural demolition. Engineering controls for stability (sequence, temporary works, exclusion zones). Administrative controls (RAMS, permit-to-work, toolbox talks, monitoring). PPE as last resort (RPE FFP3 or higher for dust; hearing protection above 85 dB(A)). \\
Cadence & Review on discovery of unexpected conditions, ACMs, weather events, near-miss or incident, change in scope or programme, or CDM notification threshold. \\
Assurance & Use audit, incident learning, and governance reporting to verify effectiveness. Structural engineer sign-off at key stages. UKAS-accredited asbestos analysis. CDM F10 notification for notifiable projects. \\
\bottomrule
\end{tabularx}
\end{table}

\begin{itemize}
\bitem{Key takeaway 1}{A pre-demolition structural survey and R\&D asbestos survey are non-negotiable prerequisites for any demolition or intrusive refurbishment project. Risk assessments produced without adequate survey data are fundamentally flawed and create legal as well as safety exposure.}
\bitem{Key takeaway 2}{The demolition sequence is the primary engineering control for structural stability. BS 6187:2011 requires that elements are removed in the reverse order of their installation. Any departure from the approved sequence must be authorised by a structural engineer and documented in a revised RAMS.}
\bitem{Key takeaway 3}{Connected services must be isolated at source and confirmed in writing before demolition commences. Visual inspection of service plans is not sufficient; trial excavation or probing should be used to verify positions where uncertainty exists.}
\bitem{Key takeaway 4}{Exclusion zones sized to encompass the collapse envelope of the structure must be physically established and enforced by barriers and banksman attendance. The collapse envelope for multi-storey structures significantly exceeds the building footprint.}
\bitem{Key takeaway 5}{The demolition risk assessment is a living document that must be reviewed and updated whenever structural conditions change, previously unidentified ACMs are discovered, significant weather events occur, or any near-miss or incident arises. Dangerous occurrences must be reported under RIDDOR 2013 and will trigger HSE investigation.}
\end{itemize}

\barquote{In Chapter~\ref{chap:Tunnelling and Underground Works Risk Assessment}, the equally demanding environment of tunnelling and underground works is examined: a discipline where the ground itself is both the medium of construction and the source of its most unpredictable hazards.}
